\documentclass[11pt]{article}
\usepackage[a4paper, top=18mm, bottom=18mm, left=20mm, right=20mm]{geometry}
\usepackage[T2A]{fontenc}
\usepackage[utf8]{inputenc}
\usepackage[ukrainian]{babel}
\usepackage{graphicx}
\setlength{\parindent}{0pt}
\setlength{\parskip}{4pt}
\pagestyle{empty}
\begin{document}
%begin
{\large\textbf{Контрольна робота}}

\textbf{Варіант \No\,1}\hfill \textit{ПІБ:}\,\underline{\hspace{60mm}}
\vspace{4mm}

%beginitem
1. Що буде виведено за виконання фрагменту коду?

%code
try:
    t = 5
    while t<=7:
        t += 2
    print(t, end=';')
except:
    print('Z')

%code

%enditem
%beginitem
2. %goodpath:Scripts/2018/Mod2018_1/03-good.txt
Відмітити все, що є коректним оператором С++:
( 1 ) cin<<3;

( 2 ) return *;

( 3 ) x=y=z

( 4 ) cout<<w;

%enditem
%beginitem
3. Що буде виведено за виконання фрагменту коду?
Записати ВИРАЗ, значенням якого є множина, що складається з кубівцілих чисел від 12 до 2002 (включно), що не діляться на жодне з чисел 2, 3.
%enditem
%beginitem
4. Що буде виведено за виконання фрагменту коду?

a,b,c=5,9,0
def h(b):
global c    a=3
    b=4
    c=5
    return a+b+c

a,b,c=6,8,7
print(h(b),a,b,c)

%enditem
%beginitem
5. Що буде виведено за виконання фрагменту коду?

\begin{center}
	\adjustimage{max size={0.8\linewidth}{0.8\paperheight}}
	{../BinTree/trees_exp/65.png}
\end{center}
Указати послідовність міток вершин дерева, зображеного на рис., що утвориться в результаті обходу дерева в оберненому порядку. Мітки записувати через пробіл.\par Обхід дерева починається з кореня (на рис. це верхня вершина).
%answer:65 оберненому

%enditem
%beginitem
6. Що буде виведено за виконання фрагменту коду?
Записати ВИРАЗ, значенням якого є список, що складається з кубівцілих чисел від 12 до 2002 (включно), що не діляться на жодне з чисел 2, 3, записаних в обернено\-му порядку.

%enditem
%beginitem
7. Що буде виведено за виконання фрагменту коду?
def h(a,b):
      c=75
      if b!=-3: c=5
      elif a>-2: return 9
      else: return 0
      return c
   print(h(-7,-7))
%enditem
%beginitem
8. Що буде виведено за виконання фрагменту коду?

%code
f=['0','1','2','3']
del f[:6]
print(f)
%code


%enditem
%beginitem
9. Що буде виведено за виконання фрагменту коду?
Задано тип вузла списку
struct Node \{int n; Node *prev, *next;\};
У списку зберігаються послідовні цілі числа від -31 до 45. Вказівник p встановлено на вузол, в якому зберігається число -7.
Чому дорівнює значення p->next->next->next->next->next->n ?
%enditem
%beginitem
10. Що буде виведено за виконання фрагменту коду?

%code
print(False > 7 >= 4 == 7.0)
%code

%enditem
%beginitem
11. Що буде виведено за виконання фрагменту коду?

print('R', end=' ')
t=7
while t>=5:
    t+=-2
print(t)


%enditem
%beginitem
12. Що буде виведено за виконання фрагменту коду?

a,b,c=9,8,6
def h(a,b=7,c=9):
    print(a,b,c,end="")

h(5,c=2,b=0)
print(a,b,c)


%enditem
%beginitem
13. Що буде виведено за виконання фрагменту коду?

%code
a,b,c=5,9,0
def h(b):
global c    a=3
    b=4
    c=5
    return a+b+c

a,b,c=6,8,7
print(h(b),a,b,c)
%code

%enditem
%beginitem
14. Що буде виведено за виконання фрагменту коду?
Написати програму, що запитує у користувача значення дійсних параметрів $x$ та $y$, обчислює для них значення виразу 
$$\frac{\log_{2} x - 53}{(\tg x - 0.3) (y - 42)}$$ 
та повiдомляє введенi данi та результат обчислення.


%enditem
%beginitem
15. Що буде виведено за виконання фрагменту коду?
a = 5
b = 9
c = 0

def h(b):
global c    a *= 3
    b = 4
    c = 5
    return a + b + c

a = 6
b = 8
c = 7

try:
    print(h(b), a, b, c, sep=';')
except:
    print('Z', a, b, c, sep=';')

%enditem
%beginitem
16. Що буде виведено за виконання фрагменту коду?
%code

def f(n):
    print("f", end="");
    return n!=-2

print(f(-7) or f(-6))
%code

%enditem
%beginitem
17. Що буде виведено за виконання фрагменту коду?
try:
    t = 5
    while t<=7:
        t += 2
    print(t, end=';')
except:
    print('Z')

%enditem
%beginitem
18. Що буде виведено за виконання фрагменту коду?
"False"!=True
%enditem
%beginitem
19. Що буде виведено за виконання фрагменту коду?
notFalse!=7 or not4>7.0 or 2<4.0
%enditem
%beginitem
20. Що буде виведено за виконання фрагменту коду?

%code
f=['0','1','2','3']
print(f.pop(-3), end=' ')
print(f)
%code


%enditem
%beginitem
21. Що буде виведено за виконання фрагменту коду?

print('R', end=' ')
f=[0,1,2,3,4,5,6,7,8,9]
print(f[:-13:-3])


%enditem
%beginitem
22. Що буде виведено за виконання фрагменту коду?

%code
#include <iostream>

int h(int a, int b){
    int c = 75;
    if (b != -3)
        c = 5;
    else if (a > -2)
         return 9;
    else 
        return 0;
    return c;
}

int main(){
    std::cout << h(-7, -7);
    return 0;
}
%code

%enditem
%beginitem
23. Що буде виведено за виконання фрагменту коду?
%code
for f in range(-13,-13,-3):
    if f>=8:
        break
        print(f,end=' ')
    if f>=8:
        break
else:
    print(f, end=' ')
print(f, end=' ')
%code

%enditem
%beginitem
24. Що буде виведено за виконання фрагменту коду?
Записати ВИРАЗ, значенням якого є список, що складається з цілих чисел від 12 до 2002 (включно), причому спочатку за спаданням записано всі, що діляться на 7, а потім решту за спаданням.
%enditem
%beginitem
25. Що буде виведено за виконання фрагменту коду?

\begin{center}
	\adjustimage{max size={0.8\linewidth}{0.8\paperheight}}
	{../BinTree/trees_exp/65.png}
\end{center}
Указати послідовність міток вершин дерева, зображеного на рис., що утвориться в результаті обходу дерева в оберненому порядку. Мітки записувати через пробіл.\par Алгоритм починає роботу з кореня (на рис. це верхня вершина).
%answer:65 оберненому

%enditem
%beginitem
26. Що буде виведено за виконання фрагменту коду?
%code
#include <iostream>
using namespace std;

int f(int &x, int &y){
    x = 3;
    y-= 4;
    return y;
}

int main(){
    int a = 9, b = 1;
    b = f(b, b);
    cout << a << ":" << b <<':';
    {
        int a = 8, b = 7;
        cout << ((b>=5) || ((a-=2) >= 5)) << ':';
        cout << a << ":" << b << ':';
    }
    {
        inta = 2;
        cout << a << ':';
    }
    cout << a << ":" << b << endl;
    return 0;
}
%code


%enditem
%beginitem
27. Що буде виведено за виконання фрагменту коду?
try:
    for f in range(-13, -13, -3):
        if f>=8:
            break
        print(f, end=';')
        if f>=8:
            break
    else:
        print(f,end=';')
    print(f)
except:
    print('ERR')

%enditem
%beginitem
28. Що буде виведено за виконання фрагменту коду?
def f(n):
    print('f', end='')
    return n!=-2

try:
    print(f(-7) or f(-6))
except:
    print('ERR')

%enditem
%beginitem
29. Що буде виведено за виконання фрагменту коду?

%code
print(5//5%5*5)
%code

%enditem
%beginitem
30. Що буде виведено за виконання фрагменту коду?

%code
cout << (!false!=7 || !4> 7.0 || 2<4.0);
%code

%enditem
%beginitem
31. Що буде виведено за виконання фрагменту коду?
f=('a','b','c','d','e',0,1,2,3); print(f[-11])
%enditem
%beginitem
32. Що буде виведено за виконання фрагменту коду?
Задано тип вузла списку struct Node {int n; Node *prev, *next;};
У списку зберігаються послідовні цілі числа від -31 до 45. Вказівник p встановлено на вузол, в якому зберігається число -7.
Чому дорівнює значення p->next->next->next->next->next->n ?
%enditem
%beginitem
33. Що буде виведено за виконання фрагменту коду?

%code
a = 5
b = 9
c = 0

def h(b):
global c    a=3
    b=4
    c=5
    return a+b+c

a, b, c = 6,8,7
try:
    print(h(b), a, b, c, sep=';')
except:
    print('Z', a, b, c, sep=';')
%code

%enditem
%beginitem
34. Що буде виведено за виконання фрагменту коду?

print('R', end=' ')
t=5
while t<=7:
    t+=2
print(t)


%enditem
%beginitem
35. Що буде виведено за виконання фрагменту коду?

%code
try:
    t = 7
    while t>=5:
        t += -2
    print(t, end=';')
except:
    print('Z')

%code

%enditem
%beginitem
36. Що буде виведено за виконання фрагменту коду?

%code
cout << (5 % 5 % 5 % 5);
%code

%enditem
%beginitem
37. Що буде виведено за виконання фрагменту коду?

print('R', end=' ')
t=7
while t>=5:
    t+=-2
print(t)


%enditem
%beginitem
38. Що буде виведено за виконання фрагменту коду?

print('R', end=' ')
f=['0','1','2','3','4']
del f[6:6]
print(f)


%enditem
%beginitem
39. Що буде виведено за виконання фрагменту коду?
%code
if (15 != 15)
    cout << "p";
else
    cout << "j";
%code


%enditem
%beginitem
40. Що буде виведено за виконання фрагменту коду?
Рядки журналу через двокрапку містять ім'я користувача (ідентифікатор), час події,тип події, код події, наприклад
\begin{verbatim}
s="username : 2022-05-05 13:26 : ERROR : 404"
\end{verbatim}
у коді 
\begin{verbatim}
res = re.fullmatch('???', s) 
s = ''
code_str = ??? 
\end{verbatim}
чим треба замінити кожен з ???, щоб значенням code\_str був код події? 



%enditem
%beginitem
41. Що буде виведено за виконання фрагменту коду?

%code
a = 5
b = 9
c = 0

def h(b):
global c    a*=3
    b=4
    c=5
    return a+b+c

a, b, c = 6, 8, 7
try:
    print(h(b), a, b, c, sep=';')
except:
    print('Z', a, b, c, sep=';')
%code

%enditem
%beginitem
42. Що буде виведено за виконання фрагменту коду?
try:
    res = "False"!=True
    if isinstance(res, bool):
        print(f'bool;{res}')
    elif isinstance(res, int):
        print(f'int;{res}')
    elif isinstance(res, float):
        print(f'float;{res:.2f}')
    else:
        print('???')
except:
    print('ERR')

%enditem
%beginitem
43. %goodpath:Scripts/2019/Mod2019_1/Lex/03-good.txt
Відмітити все, що є коректним оператором С++:
( 1 ) "a"d"

( 2 ) if a<b a=0

( 3 ) x<y<z

( 4 ) print(3,2, sep="2", end='1')

%enditem
%beginitem
44. Що буде виведено за виконання фрагменту коду?
Знайти кількість відношень на n-елементній множині, які неантирефлексивні або неантисиметричні
%enditem
%beginitem
45. Що буде виведено за виконання фрагменту коду?
Скільки 2017-елементних підмножин множини натуральних від 1 до 20172017 містять принаймні одне число, що менше 172015
%enditem
%beginitem
46. Що буде виведено за виконання фрагменту коду?
try:
    print(5, end=';')
    print(9j!=5j, end=';')
    print(0, end=';')
except TypeError:
    print(6, end=';')
except ZeroDivisionError:
    print(8, end=';')
except:
    print('ERR', end=';')
else:
    print(7, end=';')
finally:
    print(1)

%enditem
%beginitem
47. Що буде виведено за виконання фрагменту коду?

%code
print(9**-1.0*False)
%code

%enditem
%beginitem
48. Що буде виведено за виконання фрагменту коду?
В умовах попередньої задачі було виконано p->next=p->next->next;

Чому дорівнює значення p->next->next->next->next->n ?
%enditem
%beginitem
49. Що буде виведено за виконання фрагменту коду?

print('R', end=' ')
t=5
while t<=7:
    t+=2
print(t)


%enditem
%beginitem
50. Що буде виведено за виконання фрагменту коду?

def f(n):
    print("f", end="");
    return n!=-2

print(f(-7) or f(-6))

%enditem
%beginitem
51. Що буде виведено за виконання фрагменту коду?
Записати ВИРАЗ, значенням якого є множина, що складається з цілих чисел від 12 до 2002 (включно), що діляться на 7.
%enditem
%beginitem
52. Що буде виведено за виконання фрагменту коду?

%code
print('R', end=' ')
f=['0','1','2','3','4']
del f[6:6]
print(f)
%code


%enditem
%beginitem
53. Що буде виведено за виконання фрагменту коду?
Рядки журналу через = містять ім'я користувача (ідентифікатор), час події,тип події, код події, наприклад
\begin{verbatim}
s="username = 2022-05-05 13:26 = ERROR = 404"
\end{verbatim}
у коді 
\begin{verbatim}
res = re.fullmatch('???', s) 
s = ''
code_str = ??? 
\end{verbatim}
чим треба замінити кожен з ???, щоб значенням code\_str був код події? 



%enditem
%beginitem
54. Що буде виведено за виконання фрагменту коду?

%code
print("False"!=True)
%code

%enditem
%beginitem
55. Що буде виведено за виконання фрагменту коду?

%code
f = ('a','b','c','d','e',0,1,2,3)
print(f[-11])
%code


%enditem
%beginitem
56. Що буде виведено за виконання фрагменту коду?

%code
#include <iostream>
using namespace std;

int a = 5, b = 9, c = 0;

int h(int &b){
int c;    a = 3;
    b *= 4;
    c = 5;
    return a + b + c;
}

int main(){
    inta = 6;
    intb = 8;
    intc = 7;
    cout << h(b) << ':';
    cout << a << ':' << b << ':' << c;
    return 0; 
}
%code

%enditem
%beginitem
57. Що буде виведено за виконання фрагменту коду?

print('R', end=' ')
f=['0','1','2','3','4']
print(f.pop(0), end=' ')
print(f)


%enditem
%beginitem
58. %goodpath:Scripts/2019/Mod2019_1/Lex/01-good.txt
Відмітити все, що є однією правильною лексемою:
( 1 ) 1+1j

( 2 ) c2,4

( 3 ) LIFO

( 4 ) true

%enditem
%beginitem
59. Що буде виведено за виконання фрагменту коду?
Змінні \kw{next} та \kw{prev} вузла списку позначають наступний та попередній за ним вузол відповідно. Починаючи з голови, у список записано послідовні цілі числа від~$-31$ до~$45$. Нехай змінна \kw{p} позначає вузол, що зберігає значення~$-7$.
Яке число зберігається у вузлі \kw{p.next.next.next.next.next} ?

%enditem
%beginitem
60. Що буде виведено за виконання фрагменту коду?

%code
print(False>7>4>7.0)
%code

%enditem
%beginitem
61. Що буде виведено за виконання фрагменту коду?

print('R', end=' ')
f=[0,1,2,3,4,5,6,7,8,9]
print(f[:-13:-3])


%enditem
%beginitem
62. Що буде виведено за виконання фрагменту коду?

%code
print('R', end=' ')
for f in range(-8,-8,3):
    if f>=8:
        break
        print(f, end=' ')
    if f>=8:
        break
    else:
        print(f, end=' ')
    print(f, end=' ')
%code

%enditem
%beginitem
63. %goodpath:Scripts/2018/Mod2018_1/01-good.txt
Відмітити все, що є однією правильною лексемою:
( 1 ) =>

( 2 ) c2,4

( 3 ) True

( 4 ) cpp

%enditem
%beginitem
64. Що буде виведено за виконання фрагменту коду?
Записати ВИРАЗ, значенням якого є множина, що складається з цілих чисел від 12 до 2002 (включно), що діляться на 7.
%enditem
%beginitem
65. Що буде виведено за виконання фрагменту коду?
for f in range(-13,-13,-3):
    if f>=8:
        break
        print(f,end=' ')
    if f>=8:
        break
else:
    print(f, end=' ')
print(f, end=' ')

%enditem
%beginitem
66. Що буде виведено за виконання фрагменту коду?

a,b,c=5,9,0
def h(b):
global c    a*=3
    b=4
    c=5
    return a+b+c

a,b,c=6,8,7
print(h(b),a,b,c)

%enditem
%beginitem
67. Що буде виведено за виконання фрагменту коду?
try:
    res = 5//5%5*5
    if isinstance(res, bool):
        print(f'bool;{res}')
    elif isinstance(res, int):
        print(f'int;{res}')
    elif isinstance(res, float):
        print(f'float;{res:.2f}')
    else:
        print('???')
except:
    print('Z')

%enditem
%beginitem
68. Що буде виведено за виконання фрагменту коду?
Записати ВИРАЗ, значенням якого є словник, ключами якого є цілі числа від 12 до 2002 (включно), що не діляться на жодне з чисел 2, 3, а ключам відповідають їх квадратні корені 

%enditem
%beginitem
69. Що буде виведено за виконання фрагменту коду?
Записати ВИРАЗ, значенням якого є список, що складається з кубівцілих чисел від 12 до 2002 (включно), що не діляться на жодне з чисел 2, 3, записаних в оберненому порядку.
%enditem
%beginitem
70. Що буде виведено за виконання фрагменту коду?

%code
for f in range(-8,-8+4,3):
    if f>=8:
        pass
    print(f, end=' ')
    f=-8
    if f>=8:
        break
else:
    print(f, end=' ')
print(f, end=' ')
%code

%enditem
%beginitem
71. Що буде виведено за виконання фрагменту коду?
5//5%5*5
%enditem
%beginitem
72. Що буде виведено за виконання фрагменту коду?
try:
    res = 9**-1.0*False
    if isinstance(res, bool):
        print(f'bool;{res}')
    elif isinstance(res, int):
        print(f'int;{res}')
    elif isinstance(res, float):
        print(f'float;{res:.2f}')
    else:
        print('???')
except:
    print('ERR')

%enditem
%beginitem
73. Що буде виведено за виконання фрагменту коду?

%code
print('R', end=' ')
f = ('a','b','c','d','e',0,1,2,3)
print(f[-11])
%code


%enditem
%beginitem
74. Що буде виведено за виконання фрагменту коду?

%code
"False"!=True
%code

%enditem
%beginitem
75. Що буде виведено за виконання фрагменту коду?

%code
#include <iostream>
using namespace std;

bool f(int n){
    cout<<"f";
    return n!=-2;
}

int main(){
    cout<<(f(-7) || f(-6));
    return 0;
}
%code

%enditem
%beginitem
76. Що буде виведено за виконання фрагменту коду?

%code
#include <iostream>

int h(int d){
    int z = 75;
    if (d != -3) 
        z = 5;
    else if (d > -2)
         return 9;
    else
         return 0;
    return z;
}

int main(){
    std::cout << h(-7);
    return 0;
}
%code


%enditem
%beginitem
77. Що буде виведено за виконання фрагменту коду?

\begin{center}
	\adjustimage{max size={0.8\linewidth}{0.8\paperheight}}
	{../BinTree/trees_exp/65.png}
\end{center}
Указати послідовність міток вершин дерева, зображеного на рис., що утвориться в результаті обходу дерева в оберненому порядку. ̳тки записувати через пробіл.\par Обхід дерева починається з кореня (на рис. це верхня вершина).
%answer:65 оберненому
%enditem
%beginitem
78. Що буде виведено за виконання фрагменту коду?

print('R', end=' ')
for f in range(-13,-13,-3):
    if f>=8:
        break
        print(f, end=' ')
    if f>=8:
        break
    else:
        print(f, end=' ')
    print(f)

%enditem
%beginitem
79. Що буде виведено за виконання фрагменту коду?
Записати ВИРАЗ, значенням якого є множина, що складається з кубівцілих чисел від 12 до 2002 (включно), що не діляться на жодне з чисел 2, 3.

%enditem
%beginitem
80. Що буде виведено за виконання фрагменту коду?
a,b,c=5,9,0
   def h(b):
     global c     a=3
     b=4
     c=5
     return a+b+c
   a,b,c=6,8,7
   print(h(b),a,b,c)
%enditem
%beginitem
81. Що буде виведено за виконання фрагменту коду?
for f in range(-13,-13,-3):
    if f>=8:
        break
        print(f,end=' ')
    if f>=8:
        break
    else:
        print(f,end=' ')
    print(f)

%enditem
%beginitem
82. Що буде виведено за виконання фрагменту коду?
%Оригинал был утерян. Требует отладки!
%code
__d = 0

class D:
    __d = 1
    
    def __init__(self):
        self.__d = 2
        __d = 3
        D.__d = 4

obj = D()
D.__d = 5
try:
    print(__d, end=' ')
    print(obj.__d, end=' ')
    print(D.__d, end=' ')
    print(obj.__d, end=' ')
    print(obj._D__d)
except:
    print('ERR')
%code

%enditem
%beginitem
83. Що буде виведено за виконання фрагменту коду?

\begin{center}
	\adjustimage{max size={0.8\linewidth}{0.8\paperheight}}
	{../15BinTree/trees_exp/65.png}
\end{center}
Указати послідовність міток вершин дерева, зображеного на рис., що утвориться в результаті обходу дерева в оберненому порядку. Мітки записувати через пробіл.\par Алгоритм починає роботу з кореня (на рис. це верхня вершина).
%answer:65 оберненому

%enditem
%beginitem
84. Що буде виведено за виконання фрагменту коду?
$a$ є цілочисловим масивом типу $numpy.ndarray$ розмірів $5{\times}5$. Сума елементів масиву $a$ дорівнює 22. Чому дорівнює сума елементів масиву $a$ після виконання 
\begin{verbatim}
a -= 22
\end{verbatim}


Якщо код некоректний, то в якості відповіді зазначити ERROR.



%enditem
%beginitem
85. Що буде виведено за виконання фрагменту коду?
f=['0','1','2','3','4']; del f[6:6]; print(f)
%enditem
%beginitem
86. Що буде виведено за виконання фрагменту коду?
Рядки журналу через = містять ім'я користувача (ідентифікатор), час події,тип події, код події, наприклад
\begin{verbatim}
s="username = 2022-05-05 13:26 = ERROR = 404"
\end{verbatim}
у коді 
\begin{verbatim}
res = re.fullmatch('???', s) 
s = ''
code_str = ??? 
\end{verbatim}
чим треба замінити кожен з ???, щоб значенням code\_str був код події? 



%enditem
%beginitem
87. Що буде виведено за виконання фрагменту коду?
9**-1.0*False
%enditem
%beginitem
88. Що буде виведено за виконання фрагменту коду?

%code
#include <iostream>
using namespace std;

int h(int a, int b){
    int c = 75;
    if (b != -3)
        c = 5;
    else if (a > -2)
         return 9;
    else 
        return 0;
    return c;
}

int main(){
    cout << h(-7, -7);
    return 0;
}
%code

%enditem
%beginitem
89. Що буде виведено за виконання фрагменту коду?
Написати програму, що запитує у користувача значення дійсних параметрів $x$ та $y$, обчислює для них значення виразу та повiдомляє введенi данi та результат обчислення.
$$\frac{\log_{2} x - 53}{(\tg x - 0.3) (y - 42)}$$ 


%enditem
%beginitem
90. Що буде виведено за виконання фрагменту коду?
Записати ВИРАЗ, значенням якого є список, що складається з цілих чисел від 12 до 2002 (включно), причому спочатку за спаданням записано всі, що діляться на 7, а потім решту за спаданням.

%enditem
%beginitem
91. Що буде виведено за виконання фрагменту коду?

%code
a,b,c=5,9,0
def h(b):
    a*=3
    b=4
    c=5
    return a+b+c

a,b,c=6,8,7
print(h(b),a,b,c)
%code

%enditem
%beginitem
92. Що буде виведено за виконання фрагменту коду?
"False"!=True
%enditem
%beginitem
93. Що буде виведено за виконання фрагменту коду?
f=['0','1','2','3','4']; del f[6:6]; print(f)
%enditem
%beginitem
94. Що буде виведено за виконання фрагменту коду?
def h(d):
      z=75
      if d!=-3: z=5
      elif d>-2: return 9
      else: return 0
      return z
   print(h(-7))
%enditem
%beginitem
95. Що буде виведено за виконання фрагменту коду?
try:
    t = 7
    while t>=5:
        t += -2
    print(t, end=';')
except:
    print('Z')

%enditem
%beginitem
96. Що буде виведено за виконання фрагменту коду?

%code
cout << (false != 7 > 4 < 7.0);
%code

%enditem
%beginitem
97. Що буде виведено за виконання фрагменту коду?
t=7
   while t>=5: t+=-2
   print(t)
%enditem
%beginitem
98. Що буде виведено за виконання фрагменту коду?

%code
cout << 5 % 5 % 5 % 5;
%code

%enditem
%beginitem
99. Що буде виведено за виконання фрагменту коду?
f=['0','1','2','3','4']; print(f.pop(0),end=' '); print(f)
%enditem
%beginitem
100. Що буде виведено за виконання фрагменту коду?
Записати ВИРАЗ, значенням якого є список, що складається з цілих чисел від 12 до 2002 (включно), причому спочатку за спаданням записано всі, що діляться на 7, а потім решту за спаданням.
%enditem




































































































%end
\newpage
\end{document}


